\chapter{Multiband Aperture and Polarimetry / Abertura Multifaixa e Polarimetria}

\section*{English}
The aperture is intentionally hybrid rather than uniform across the spectrum. One external VHF Yagi and one external UHF Yagi share a mast and provide orthogonal linear polarizations: the VHF antenna uses larger elements and the UHF antenna uses smaller elements. Smaller L/S/C, X/Ku and K/Ka antennas or tiles occupy the face layers and local mounts. HF uses loops, crossed dipoles and platform-mode diversity. Each antenna family has its own filtering, conversion and sampling strategy, while the face preserves a common timing, calibration and fusion interface.

Each relevant element provides two coherent orthogonal ports. In a local $x$--$y$ basis:
\[
E_{\mathrm{RHCP}}=\frac{E_x-jE_y}{\sqrt{2}},\qquad E_{\mathrm{LHCP}}=\frac{E_x+jE_y}{\sqrt{2}}.
\]
For a curved surface, local bases must be rotated to a common coordinate system before fusion. Polarimetric calibration is a function of frequency, angle, temperature, amplitude balance, phase and mutual coupling.

\subsection*{Engineering innovation / Inovação de engenharia}
The crossed-Yagi solution is a central engineering innovation of this project. It converts a difficult broadband and polarization requirement into a mechanically integrated dual-aperture assembly: the larger Yagi covers VHF, the smaller Yagi covers UHF, and their transverse elements are arranged in orthogonal planes for 90-degree polarization diversity. Each antenna remains independently matched, digitized and calibrated, while the face-level interface combines their observations. This provides aggregate spectral coverage and polarization diversity without claiming that one physical Yagi covers both bands.

\begin{figure}[H]
\centering
\includegraphics[width=0.95\textwidth]{fig04_particionamento_espectro.png}
\caption{Recommended spectral partition / Particionamento espectral recomendado.}
\label{fig:spectrum}
\end{figure}

\begin{figure}[H]
\centering
\includegraphics[width=0.95\textwidth]{fig05_polarimetria.png}
\caption{Vector polarimetric acquisition / Aquisição polarimétrica vetorial.}
\label{fig:polarimetry}
\end{figure}

\section*{Português}
A abertura é deliberadamente híbrida, e não uniforme em todo o espectro. Uma Yagi VHF externa e uma Yagi UHF externa compartilham o mastro e fornecem polarizações lineares ortogonais: a antena VHF usa elementos maiores e a antena UHF usa elementos menores. Antenas ou tiles menores de L/S/C, X/Ku e K/Ka ocupam as camadas da face e seus suportes locais. HF utiliza loops, dipolos cruzados e diversidade de modos da plataforma. Cada família de antenas possui filtragem, conversão e amostragem próprias, enquanto a face preserva uma interface comum de sincronismo, calibração e fusão.

Cada elemento relevante fornece duas portas ortogonais coerentes. Em uma base local $x$--$y$:
\[
E_{\mathrm{RHCP}}=\frac{E_x-jE_y}{\sqrt{2}},\qquad E_{\mathrm{LHCP}}=\frac{E_x+jE_y}{\sqrt{2}}.
\]
Em uma superfície curva, as bases locais devem ser rotacionadas para um sistema comum antes da fusão. A calibração polarimétrica depende de frequência, ângulo, temperatura, equilíbrio de amplitude, fase e acoplamento mútuo.

\subsection*{Inovação de engenharia}
A solução de Yagis cruzadas é uma inovação central de engenharia deste projeto. Ela transforma um requisito difícil de largura de banda e polarização em um conjunto dual de aberturas mecanicamente integrado: a Yagi maior cobre VHF, a Yagi menor cobre UHF e seus elementos transversais são dispostos em planos ortogonais para diversidade de polarização a 90 graus. Cada antena permanece casada, digitalizada e calibrada de forma independente, enquanto a interface da face combina as observações. Assim, obtém-se cobertura espectral agregada e diversidade de polarização sem afirmar que uma única Yagi cubra as duas bandas.
